%! Representing a Quantitative Variable with Graphs — Unit 1, Lesson 1.3 — Homework
% GRADUAL RELEASE: this is the lesson's SCORED practice and the LAST component in the packet.
% Paper homework is the DEFAULT for this course; DeltaMath is an occasional override that
% replaces this page and strikes its score cell on the cover. Started in the last eight minutes
% of the period and finished at home. The AP Practice component that precedes it stays
% assigned-but-unscored.
% THIRD CONTEXT: the notes teach on 24 quiz finishing times, the AP Practice page runs on a
% rainfall record, and these items are on fresh data again — recall is not the skill being built.
% \answerspace{H}{} reserves exactly H in the blank; the key passes the answer as the second
% argument, so the two files paginate identically by construction. Keep the macro and every
% \boxguard byte-identical in the blank and the key.
\documentclass[12pt]{article}
\usepackage{apstats-article}
\usepackage{apstats-boxes}
\newcommand{\answerspace}[2]{\par\nopagebreak\noindent\begin{minipage}[t][#1][t]{\linewidth}%
  \color{keyred}\bfseries #2\end{minipage}\par}

\begin{document}

\pageheader{Unit 1, Lesson 1.3}{Homework}
% NAMESTRIP: no \namedateperiod here — the name row lives on the cover only.

\vspace{0.15in}

% Authored BYTE-IDENTICALLY in homework/ and homework_key/ — this box is what keeps the
% blank and the key the same number of pages.
\boxguard[8]
\begin{remindbox}
\small
\textbf{This is your graded homework.} It is scored and due the first class after two study halls. You will start it in the
last minutes of the period, so ask about anything that stalls you before you leave, then finish
it on your own.
\end{remindbox}

\vspace{0.12in}

% ── The scored practice — five exercises spanning Topic 1.5 ─────────────────
\boxguard[14]
\begin{notesbox}{Reading and Choosing a Quantitative Display}
Show the arithmetic wherever a question asks for a number, and write every explanation
\emph{in the context of the problem} --- that is what earns credit on the AP exam, and it is the
standard here.
\begin{enumerate}[leftmargin=1.9em, itemsep=6pt]
  \item Label each variable \textbf{discrete} or \textbf{continuous}.

        \vspace{3pt}
        \renewcommand{\arraystretch}{1.3}
        \begin{tabularx}{\linewidth}{@{} Y Y Y Y @{}}
        \textbf{Cars in a lot} & \textbf{Height (cm)} & \textbf{Goals scored} &
        \textbf{Time to run 400 m} \\
        \blank{1.9cm} & \blank{1.9cm} & \blank{1.9cm} & \blank{1.9cm} \\
        \end{tabularx}
  \item A histogram of 60 test scores uses blocks of 10 points and shows one tall middle bar.
        Redrawn with blocks of 2 points, it shows two separate peaks. Which version is
        \emph{wrong}? Explain.
        \answerspace{3.1cm}{}
  \item In a histogram the bars touch; in the bar graph you drew last lesson they did not.
        Explain what that difference is telling you about the horizontal axis.
        \answerspace{3.1cm}{}
  \item A stem plot has stem \textbf{7} with leaves \textbf{2 4 4 9}. Write out the four values
        those represent. \blank{6.0cm}

        Why can a histogram of the same data not give you those four numbers back?
        \answerspace{2.4cm}{}
  \item \textbf{AP-style multiple choice.} A researcher has 200 reaction times measured to the nearest
        thousandth of a second and wants a display that shows the overall shape. Which
        statement is \textbf{correct}? Circle one, then justify it in a sentence.
        \begin{enumerate}[label=\textbf{(\Alph*)}, leftmargin=2.2em, itemsep=2pt]
          \item A stem plot is best, because it keeps every individual value.
          \item A histogram is reasonable, because with 200 values the shape matters more than
                the individual numbers.
          \item A bar graph is best, because reaction time is measured numerically.
          \item No display works, because the values are continuous.
          \item A histogram is wrong, because it changes appearance with the interval width.
        \end{enumerate}
        \answerspace{2.4cm}{}
\end{enumerate}
\end{notesbox}

\vspace{0.12in}

% The next-lesson preview closes the packet, so it lives here rather than in AP Practice.
\boxguard[10]
\begin{spiralbox}
\small
\textbf{Next lesson: Describing the Distribution of a Quantitative Variable.} You can now build
and read these displays. Next time we put words to what they show --- where the values pile up,
how spread out they are, whether the picture leans one way, and which values sit far from the
rest --- and you will learn the four things every AP answer about a display has to mention.
\end{spiralbox}

\end{document}
